\section{Calculus of expansions with transported precision}
\label{sec:series-calculus}

The operations \wl{SeriesAdd}, \wl{SeriesMultiply} and \wl{SeriesPower}
provide arithmetic without discarding remainders. The functions
\wl{SeriesLog} and \wl{SeriesExp} handle logarithms and exponentials;
\wl{SeriesCompose} and \wl{SeriesObservable} handle composition and other
observables. Truncation, refinement and justified differentiation use
\wl{SeriesTruncate}, \wl{SeriesRefine} and \wl{SeriesDifferentiate}.
All ten operations share a common representation
\begin{equation}
 F(y)=b(y)+c(y)\left[J(w(y))+\OO(w(y)^P\M(w(y))^D)\right],
 \qquad w(y)\longrightarrow0^+,
 \label{eq:calculus-carrier}
\end{equation}
where $b,c,w$ are exact expressions and $J$ is a finite power--log jet.
The offset $b$ and prefactor $c$ are recorded separately. Thus a cutoff in
the jet describes relative precision with respect to $c$, and the absolute
remainder contains $|c|$. Ordinary forward and inverse objects, Lambert
objects, and explicitly recorded coordinate substitutions can be converted
to this representation. An operation returns a derived object and its replay
recipe; it does not retain an inverse-model label that would falsely suggest
that the original inverse residual checker applies to an arbitrary observable.

For addition and multiplication, the precision rules of
Section~\ref{sec:scale} are used after expressing both operands in the same
positive coordinate. An exact prefactor may be absorbed into the jet when
it is a finite power--log expression. Distinct prefactors can be aligned
when their ratio lies in that algebra. The package rejects incompatible
coordinates or independently ordered exponential carriers rather than
assigning them an unsupported single cutoff.

\paragraph{Powers and logarithms.}
For a positive monomial leading term $a w^\alpha$, write
$J=a w^\alpha(1+U)$. A constant power and a logarithm reduce to the unit
series of $(1+U)^r$ and $\log(1+U)$. The input uncertainty first becomes
$\OO(w^{P-\alpha}\M^D)$ in $U$. Multiplication by the output prefactor
then transports it to the correct absolute scale. Real noninteger powers
require positivity on the selected branch. These operations retain any
earlier input precision ceiling; requesting a higher output cutoff does
not create unknown coefficients.

\paragraph{Exponentials and exact carriers.}
Suppose
\[
 H=H_{\le0}+H_{>0}+\OO(w^P\M^D),\qquad P>0,
\]
where $H_{\le0}$ contains the nonpositive-weight blocks. Then
\[
 e^H=e^{H_{\le0}}
 \left[e^{H_{>0}}+\OO(w^P\M^D)\right].
\]
Indeed the omitted argument tends to zero and $e^r=1+\OO(r)$ there.
The operation keeps $e^{H_{\le0}}$ as an exact carrier and expands only
the vanishing part. For example $H=w^{-1}+\log w+w$ yields
$w e^{1/w}(1+w+w^2/2+\cdots)$. This also supplies the precision rule
needed when a source logarithm is inverted and the source variable must
be reconstructed by exponentiation. A nonvanishing argument uncertainty
is rejected: it cannot justify a relative error tending to zero after
exponentiation.

\paragraph{Composition and regular observables.}
The outer local coordinate must become $a w^\alpha(1+o(1))$ with
$a,\alpha>0$ after substitution of the inner expansion. This condition
checks both the limiting point and its selected side. An outer remainder
$\OO(v^P\M(v)^D)$ becomes $\OO(w^{\alpha P}\M(w)^D)$, since
$\log v=\alpha\log w+\OO(1)$. The uncertainty of the inner argument is
propagated independently through every jet operation, and the worse
precision is retained. Regular analytic observables at a finite limiting
argument use their Taylor series and the same precision arithmetic.
Poles, branch points and logarithmic leading coefficient functions require
their respective scale normalizations; a general analytic-observable call
does not silently assume them to be regular.

\paragraph{Differentiation requires derivative information.}
An estimate $R(w)=\OO(w^P\M(w)^D)$ alone does not imply any estimate on
$R'$. For example $w^P\sin(e^{1/w})$ has that magnitude bound and a much
larger derivative. Therefore a nonzero remainder can be differentiated
$n$ times only with a recorded or explicitly supplied contract
\[
 R^{(j)}(w)=\OO(w^{P-j}\M(w)^D),\qquad 0\le j\le n.
\]
The option \wl{"RemainderDerivativeOrder" -> n} is a declaration of
these hypotheses, not a numerical test of them. Each differentiation
consumes one order of the contract. The jet derivative uses
$D_w[w^\beta Q(\log w)]=w^{\beta-1}(\beta Q+Q')$; the derivatives
of the exact coordinate, offset and prefactor in
\eqref{eq:calculus-carrier} are included by the chain and product rules.
Exact finite expressions require no remainder hypothesis. A discarded tail
of a known finite power--log expression also has derivative bounds of every
order. In contrast, a derivative bound declared for one unknown remainder
is not automatically promoted to a stronger exponent when the source is
refined; that stronger bound must be established separately.

\paragraph{Truncation and refinement.}
Truncation moves the first discarded block, including its logarithmic
degree, into the remainder. It cannot improve uncertainty inherited from
the input. Refinement instead replays a retained forward expression,
inverse problem, coordinate substitution or operation recipe at a higher
working order. The replay still respects transported input precision
ceilings. Without a retained source or recipe, the package reports that
refinement is unavailable; changing an inert remainder exponent would not
be a mathematical refinement.

An automatically computed forward remainder is distinct from a declared
input error. The object records the latter separately as
\wl{DeclaredInputRemainder}. When refining an analytic input such as
$\sin x$, the original function can be expanded further; its previous
Taylor remainder must not become a permanent declared accuracy ceiling.
A remainder explicitly supplied for an unspecified omitted function
remains a binding ceiling. Tests exercise both cases with independent
inverse coefficients.
